% =========================================================
% APPENDIX A: Laplacian conditioning
% =========================================================
\chapter{Laplacian conditioning}
\label{app:laplacian}

\section{Residual training for Schr\"odinger}
\label{app:laplacian:setup}

We consider the stationary many-body Schr\"odinger equation on configuration space
$\mathbf R=(\mathbf r_1,\ldots,\mathbf r_N)\in\Omega\subset\mathbb R^{D}$ with $D=dN$
(for our quantum dots, $d=2$):
\begin{equation}
\hat H\Psi(\mathbf R)=E\Psi(\mathbf R),
\qquad
\hat H=\hat T+\hat V,
\qquad
\hat T=-\frac12\sum_{i=1}^N\Delta_i,
\label{eq:lap:sch}
\end{equation}
where $\hat V$ collects external and interaction potentials.

In \emph{strong residual minimisation} one introduces
\begin{equation}
L := \hat H - E \quad\Rightarrow\quad L\Psi = 0,
\end{equation}
and minimises a squared residual loss of the schematic form
\begin{equation}
\mathcal L_{\mathrm{res}}(\theta)
= \frac12\|L\Psi_\theta\|_{L^2(\Omega;w)}^2
= \frac12\int_\Omega |(L\Psi_\theta)(\mathbf R)|^2\,w(\mathbf R)\,d\mathbf R,
\label{eq:lap:resloss}
\end{equation}
for a nonnegative weight $w$ induced by the sampling strategy.
A key message of the numerical-analysis viewpoint is that strong residual losses
can be \emph{severely ill-conditioned}, and training speed is tightly linked to
spectral properties of an operator/matrix induced by $L$ \cite{De_Ryck_2024}.

\paragraph{Which of the thesis's objectives this analyses.} Three objectives appear in the
main text and Eq.~\eqref{eq:lap:resloss} is exactly one of them. Stage~I pretraining
(\S\ref{sec:residual-stage}) minimises $\frac1B\sum_b (E_L(\tilde R_b)-E_{\rm eff})^2$
over configurations drawn from $|\Psi_\theta|^2$; since $E_L=\hat H\Psi/\Psi$, that sample
average estimates $\int |\Psi_\theta|^2\,(\hat H\Psi_\theta/\Psi_\theta - E)^2
=\|(\hat H-E)\Psi_\theta\|^2_{L^2}$, which is Eq.~\eqref{eq:lap:resloss} with $w\equiv1$.
Everything in this appendix therefore applies to stage~I directly. It does \emph{not} apply
in the same way to the other two. Stage~II is variational energy minimisation
(\S\ref{sec:sr-stage}), whose objective is linear rather than quadratic in the residual and
whose conditioning is governed by the quantum geometric tensor measured in
\S\ref{sec:kernel-cond}, not by the $\Delta^2$ stiffness derived below. The collocation
route (\S\ref{sec:method-collocation}) replaces the strong residual with a weak-form
score-function objective under an analytic proposal, for the reason developed in
\S\ref{app:laplacian:barron} and Appendix~\ref{app:catch22}: the weak form is what removes the
second-derivative amplification, and it is the escape from the mechanism analysed here
rather than an instance of it. Where the two production objectives inherit anything from
this analysis, it is through the single channel they share with stage~I---the Laplacian
inside $E_L$---and not through the squared-residual structure.

\section{Linearized geometry and the matrix \(A\)}
\label{app:laplacian:A}

To expose the mechanism, follow the standard linearization argument used in
conditioning analyses \cite{De_Ryck_2024}.
Fix an expansion point $\theta_0$ and define tangent (feature) functions
\begin{equation}
q_j(\mathbf R) := \partial_{\theta_j}\Psi_\theta(\mathbf R)\big|_{\theta=\theta_0},
\qquad j=1,\ldots,P.
\end{equation}
For small parameter increments $\delta\theta$, we have the first-order model
\begin{equation}
\Psi_{\theta_0+\delta\theta}(\mathbf R)
\approx
\Psi_0(\mathbf R)+\sum_{j=1}^P q_j(\mathbf R)\,\delta\theta_j,
\qquad
\Psi_0:=\Psi_{\theta_0}.
\label{eq:lap:linPsi}
\end{equation}
Applying the residual operator gives
\begin{equation}
(L\Psi_{\theta_0+\delta\theta})(\mathbf R)
\approx
(L\Psi_0)(\mathbf R)+\sum_{j=1}^P (Lq_j)(\mathbf R)\,\delta\theta_j.
\label{eq:lap:linRes}
\end{equation}

Introduce the weighted inner product (complex-valued wavefunctions allowed)
\begin{equation}
\langle f,g\rangle_w := \int_\Omega f(\mathbf R)^\ast\,g(\mathbf R)\,w(\mathbf R)\,d\mathbf R.
\end{equation}
Let the Jacobian $J$ be the linear map whose $j$th column is $Lq_j$.
Then the loss \eqref{eq:lap:resloss} becomes (to second order in $\delta\theta$)
\begin{equation}
\mathcal L_{\mathrm{res}}(\theta_0+\delta\theta)
\approx
\frac12\|r_0 + J\,\delta\theta\|_w^2,
\qquad
r_0 := L\Psi_0.
\label{eq:lap:lsq}
\end{equation}
with normal-equation matrix
\begin{equation}
A := J^\ast J,
\qquad
A_{ij}=\langle Lq_i,\,Lq_j\rangle_w.
\label{eq:lap:Adef}
\end{equation}
Here $A$ is Hermitian positive semidefinite; in practice it can be close to singular,
so one often effectively works with a damped variant $A+\lambda I$ (implicit or explicit).

Now consider gradient descent on the quadratic model \eqref{eq:lap:lsq}
(with step size $\eta$). The error in parameter space evolves as
\begin{equation}
\delta\theta^{(t+1)}-\delta\theta^\star
=
\left(I-\eta A\right)\left(\delta\theta^{(t)}-\delta\theta^\star\right),
\label{eq:lap:GD}
\end{equation}
so each eigen-direction of $A$ contracts by a factor $(1-\eta\lambda)$.
Hence convergence speed is governed by the eigenvalue spread of $A$,
typically summarized by the (effective) condition number
\begin{equation}
\kappa(A)=\frac{\lambda_{\max}(A)}{\lambda_{\min}(A)}
\end{equation}
(over the nonzero/regularised spectrum).
Large $\kappa(A)$ implies stiffness: a step size small enough to be stable for
$\lambda_{\max}$ yields slow progress along directions associated with
$\lambda_{\min}$. This is precisely the conditioning bottleneck emphasized in
physics-informed training analyses \cite{De_Ryck_2024}.

\section{Bi-Laplacian scaling}
\label{app:laplacian:bilap}

The defining feature of Schr\"odinger is that $L$ contains a second-order
differential operator. To isolate the kinetic contribution, consider the regime
where the residual is dominated by the Laplacian part,
\begin{equation}
L \approx -\frac12\Delta_{\mathbf R}.
\label{eq:lap:Lapprox}
\end{equation}

\paragraph{From \(L\) to \(L^\ast L\).}
The matrix $A$ in \eqref{eq:lap:Adef} is a Gram matrix in the feature space
$\{Lq_j\}$.
The corresponding operator-level object is the \emph{Hermitian square}
$L^\ast L$ (with adjoint taken in the $L^2(\Omega;w)$ inner product)
\cite{De_Ryck_2024}. For general non-constant $w$, the adjoint $L^\ast$ differs from the
formal differential adjoint by weight-induced lower-order terms; the cleanest picture
is obtained for constant $w$ and standard boundary conditions (periodic/Dirichlet),
where the Laplacian is self-adjoint.

In the Laplacian-dominant regime \eqref{eq:lap:Lapprox}, one has (up to a constant factor)
\begin{equation}
L \approx -\tfrac12\Delta
\quad\Rightarrow\quad
L^\ast L \approx \tfrac14\,\Delta^2,
\label{eq:lap:bilap}
\end{equation}
i.e.\ a bi-Laplacian scaling. The prefactor does not affect conditioning ratios.

\paragraph{Fourier/eigenmode derivation of \(k^4\).}
On a periodic domain (or, more generally, in a Laplacian eigenbasis)
\(
-\Delta \phi_{\mathbf k} = |\mathbf k|^2\phi_{\mathbf k},
\)
so
\(
\Delta^2\phi_{\mathbf k} = |\mathbf k|^4\phi_{\mathbf k}.
\)
If an error $\delta\Psi$ expands as $\delta\Psi=\sum_{\mathbf k}c_{\mathbf k}\phi_{\mathbf k}$
and $L\approx -\tfrac12\Delta$, then
\begin{equation}
\|L\,\delta\Psi\|_{L^2}^2
\approx
\frac14\sum_{\mathbf k} |\mathbf k|^4|c_{\mathbf k}|^2.
\label{eq:lap:k4}
\end{equation}
Thus the squared residual penalizes high-frequency components with a \emph{quartic}
weight $|\mathbf k|^4$.
If the relevant spectrum spans $|\mathbf k|\in[k_{\min},k_{\max}]$, the induced
eigenvalue ratio scales like
\begin{equation}
\kappa \sim \left(\frac{k_{\max}}{k_{\min}}\right)^4,
\label{eq:lap:kappa}
\end{equation}
which is a sharp mathematical expression of ``Laplacian-driven stiffness''.
This quartic growth is consistent with the Laplacian/Poisson toy examples and
frequency-dependent conditioning effects emphasized in \cite{De_Ryck_2024}.

\paragraph{Length-scale interpretation.}
Let $\ell$ denote the smallest length scale that must be represented.
Then $k_{\max}\sim 1/\ell$ and \eqref{eq:lap:kappa} implies
\begin{equation}
\kappa \ \text{grows rapidly as}\ \ell\downarrow 0.
\end{equation}
In many-body Schr\"odinger problems, small $\ell$ arises from nodal structure,
strong localisation, and short-range correlation features; in all such regimes,
the Laplacian makes the strong residual extremely curvature-sensitive.

\section{Barron/spectral perspective: the Laplacian}
\label{app:laplacian:barron}

The conditioning story above can be rephrased in approximation terms.
In spectral Barron-type viewpoints~\cite{Barron1993-UniversalApprox}, function classes are controlled by norms that
weight the Fourier transform by powers of frequency (or related spectral weights).
A key implication is that applying derivatives systematically increases the
importance of high-frequency tails \cite{De_Ryck_2024}.
Since $\Delta$ multiplies a Fourier mode by $|\mathbf k|^2$, and the squared
residual effectively involves $L^\ast L$ (hence $\Delta^2$ in Laplacian-dominant
regimes), the residual can require substantially more ``spectral budget'' than
the wavefunction itself.
In practical terms: even if $\Psi$ is approximable in a given network class,
\emph{approximating its strong residual} can be much harder, because the loss
implicitly demands accurate high-frequency curvature (cf.\ \eqref{eq:lap:k4}).
This is one of the motivations, in numerical analysis of PINNs, for considering
weak/variational formulations that reduce derivative order when the strong form
is too stiff \cite{De_Ryck_2024}.

\section{Implications for Stochastic Reconfiguration}
\label{app:laplacian:sr}

Stochastic reconfiguration (SR) updates parameters using a natural-gradient-like
preconditioner \cite{BeccaSorella2017-QMCBook,Amari1998-NaturalGradient,ParkKastoryano2020-GeometryNQS}.
Writing $O_i=\partial_{\theta_i}\log\Psi_\theta$, one common SR update is
\begin{equation}
S(\theta)\,\delta\theta = -g(\theta),
\qquad
S_{ij}=\mathrm{Cov}_{p_\theta}[O_i,O_j],
\qquad
g_i=2\,\mathrm{Cov}_{p_\theta}[E_L,O_i],
\label{eq:lap:sr}
\end{equation}
where $p_\theta\propto|\Psi_\theta|^2$ and $E_L=(\hat H\Psi_\theta)/\Psi_\theta$.
SR improves conditioning in parameter space by using the (quantum) information
geometry of the variational family \cite{Amari1998-NaturalGradient,ParkKastoryano2020-GeometryNQS}.
However, the Laplacian-induced stiffness described above is an \emph{operator-level}
phenomenon: strong residuals magnify curvature errors via $\Delta$ and $\Delta^2$
(cf.\ \eqref{eq:lap:k4}), which remains a fundamental challenge even when parameter
updates are geometrically preconditioned.



% =========================================================
% APPENDIX B: Coulomb conditioning (2D, spin channels)
% =========================================================
\chapter{Coulomb conditioning}
\label{app:coulomb}

\section{Coulomb structure and singularity}
\label{app:coulomb:setup}

For our two-dimensional electronic systems ($\mathbf r_i\in\mathbb R^2$), the interaction
contains Coulomb terms of the form
\begin{equation}
\hat V_{\mathrm C}(\mathbf R)=\sum_{1\le i<j\le N}\frac{1}{r_{ij}},
\qquad r_{ij}=|\mathbf r_i-\mathbf r_j|.
\label{eq:coulomb:V}
\end{equation}
This is a \emph{singular multiplication operator}: it diverges at particle
coalescence $r_{ij}\to 0$.
The exact eigenfunction remains finite (or vanishes by antisymmetry) but develops
non-smooth short-distance structure (a cusp) so that divergences in kinetic and
potential contributions cancel in physically meaningful quantities such as the
local energy \cite{BeccaSorella2017-QMCBook}.

\begin{quote}
Coulomb does not merely add ``another term'' to the residual; it introduces
\emph{singular short-range physics} that forces delicate cancellations with the
Laplacian. Strong residual objectives square these effects and become extremely
sensitive to tiny cusp errors.
\end{quote}

\section{2D cusp conditions: antiparallel vs parallel spins}
\label{app:coulomb:cusp}

We present the 2D cusp logic in a way that separates the two spin channels.
The cleanest statement (especially for practical ans\"atze) is in terms of the
short-distance behaviour of a \emph{Slater--Jastrow} factorization
\begin{equation}
\Psi(\mathbf R)=D(\mathbf R)\,J(\mathbf R),
\qquad
J(\mathbf R)=\exp\!\Big(\sum_{i<j} u(r_{ij})\Big),
\label{eq:coulomb:slater-jastrow}
\end{equation}
because (i) the Slater part $D$ enforces antisymmetry and determines whether $\Psi$
vanishes at coalescence, and (ii) the Jastrow slope $u'(0)$ is what cancels the
$1/r$ Coulomb singularity in the local energy \cite{BeccaSorella2017-QMCBook}.

Fix a pair $(i,j)$ and introduce relative/centre-of-mass coordinates in 2D:
\begin{equation}
\mathbf R_{ij} := \frac{\mathbf r_i+\mathbf r_j}{2},
\qquad
\mathbf r := \mathbf r_i-\mathbf r_j,
\qquad
r:=|\mathbf r|.
\end{equation}
A direct computation gives, in any dimension,
\begin{equation}
\Delta_i+\Delta_j = \frac12\,\Delta_{\mathbf R_{ij}} + 2\,\Delta_{\mathbf r},
\label{eq:coulomb:lap_decomp}
\end{equation}
hence the singular short-distance kinetic behaviour is controlled by the relative term
$-\Delta_{\mathbf r}$.

A key 2D identity is
\begin{equation}
\Delta_{\mathbf r} r = \frac{1}{r}\qquad (r>0,\ \mathbf r\in\mathbb R^2).
\label{eq:coulomb:lap_r_2d}
\end{equation}

\paragraph{Antiparallel spins (distinguishable in the Slater part).}
If the two electrons have opposite spin, they enter different determinants in $D$
and generically $D$ does \emph{not} vanish at $\mathbf r\to 0$. Then $\Psi(0)\neq 0$ and
a direct cusp expansion in 2D is legitimate:
\begin{equation}
\Psi(\mathbf r,\ldots)
=
\Psi(0,\ldots)\Big(1 + a\,r + \mathcal O(r^2)\Big),
\qquad r\to 0.
\label{eq:coulomb:expansion_anti}
\end{equation}
Using \eqref{eq:coulomb:lap_r_2d}, the singular part of the relative Laplacian is
\begin{equation}
\Delta_{\mathbf r}\Psi \approx \Psi(0)\,a\,\Delta_{\mathbf r}r
= \Psi(0)\,\frac{a}{r},
\qquad
\frac{\Delta_{\mathbf r}\Psi}{\Psi}\approx \frac{a}{r}.
\label{eq:coulomb:lap_sing_2d}
\end{equation}
Since the pair kinetic contribution contains $-\Delta_{\mathbf r}$, the singular part
of the local energy is
\begin{equation}
E_L(\mathbf R)=\frac{\hat H\Psi}{\Psi}\supset \left(-\frac{a}{r}+\frac{1}{r}\right)
=\frac{1-a}{r}.
\label{eq:coulomb:EL_sing_2d_anti}
\end{equation}
Thus finiteness of $E_L$ at coalescence requires
\begin{equation}
a=1,
\qquad\text{i.e.}\qquad
\left.\partial_r \log \Psi\right|_{r=0}=1
\quad\text{(2D, antiparallel spins)}.
\label{eq:coulomb:cusp_2d_anti}
\end{equation}
Equivalently, in a Slater--Jastrow form with $D(0)\neq 0$,
this is the same as $\partial_r\log J|_{0}=1$ for the $ij$-pair.

\paragraph{Parallel spins (identical fermions, node at coalescence).}
For same-spin electrons, antisymmetry forces $D$ (and hence $\Psi$) to vanish as
$\mathbf r\to 0$. Generically one has a \emph{linear} zero:
\begin{equation}
D(\mathbf r,\ldots) \;=\; \bm\eta(\ldots)\cdot \mathbf r \;+\; \mathcal O(r^2),
\qquad r\to 0,
\label{eq:coulomb:D_linear_zero}
\end{equation}
with some nonzero vector $\bm\eta$ depending on the other coordinates.
In this case, the cancellation of the Coulomb singularity is controlled by the
Jastrow factor through cross terms in $-\Delta_{\mathbf r}(DJ)/(DJ)$.
A standard short-distance analysis (carried out explicitly for 2D quantum dots)
shows that the divergent kinetic contribution behaves like
\begin{equation}
-\frac{\Delta_{\mathbf r}(DJ)}{DJ}
\supset
-\frac{(2k+d-1)}{r}\,\partial_r\log J
\qquad\text{with}\qquad
d=2,\ \ k=1\ \text{(identical particles)},
\label{eq:coulomb:jastrow_div}
\end{equation}
while the Coulomb potential contributes $+1/r$ to $E_L$ as before.
Hence finiteness of the local energy requires
\begin{equation}
\left.\partial_r\log J\right|_{r=0}=\frac{1}{2k+d-1}=\frac{1}{3}
\quad\text{(2D, parallel spins)}.
\label{eq:coulomb:cusp_2d_par}
\end{equation}
This $1$ (antiparallel) versus $1/3$ (parallel) 2D cusp statement is the
standard quantum-dot/Jastrow--Slater cusp condition \cite{Siljamaki2003-QD-Cusps}.

\paragraph{What to remember.}
In 2D:
\begin{align}
\text{antiparallel spins:}\quad & \left.\partial_r\log \Psi\right|_{0}=1
\;\;\;\Longleftrightarrow\;\;\;
\left.\partial_r\log J\right|_{0}=1,
\label{eq:coulomb:summary_anti}
\\
\text{parallel spins:}\quad & \Psi(0)=0\ \text{(linear node)},\ \ 
\left.\partial_r\log J\right|_{0}=\frac13.
\label{eq:coulomb:summary_par}
\end{align}

\section{Why Coulomb is ill-conditioned (2D emphasis)}
\label{app:coulomb:residual}

The cusp conditions above show \emph{how} Coulomb creates stiffness: it forces
kinetic and potential contributions to be individually large and singular,
while their \emph{sum} is finite only under a short-distance constraint.
Strong residual minimisation is therefore a \emph{cancellation problem}.

\paragraph{Sensitivity amplification in the local energy.}
If the learned antiparallel-spin cusp slope is $a_\theta = 1+\delta a$, then
\eqref{eq:coulomb:EL_sing_2d_anti} implies
\begin{equation}
E_L(\mathbf R)\supset -\frac{\delta a}{r},
\label{eq:coulomb:EL_mismatch_2d_anti}
\end{equation}
so an arbitrarily small cusp mismatch produces an arbitrarily large local-energy spike
as $r\to 0$.

For parallel spins, if the learned Jastrow slope is
$\partial_r\log J|_{0}=\frac13+\delta b$, then the same short-distance analysis leading to
\eqref{eq:coulomb:cusp_2d_par} implies
\begin{equation}
E_L(\mathbf R)\supset -\frac{3\,\delta b}{r}.
\label{eq:coulomb:EL_mismatch_2d_par}
\end{equation}
Thus both channels generate $1/r$ spikes unless the cusp is satisfied.

\paragraph{Why squaring is especially dangerous in 2D.}
A $1/r$ defect in $E_L$ (or in $(\hat H-E)\Psi$) is concentrated near coalescence.
But in 2D the squared singularity is not integrable---it diverges logarithmically:
\[
\int_{r<\varepsilon}\Big(\frac{1}{r}\Big)^2\, d^2\mathbf r
=
\int_0^\varepsilon \frac{1}{r^2}\,(2\pi r\,dr)
=
2\pi\int_0^\varepsilon \frac{dr}{r}
=+\infty.
\]
Therefore, without effectively resolving the cusp, \emph{strong} squared-residual
objectives can be dominated (or even made ill-defined in the continuum limit) by the
coalescence region. In Monte Carlo practice, this manifests as heavy-tailed fluctuations
and extreme sensitivity to rare close-pair samples.

\paragraph{Residual expansion and cancellation structure.}
Writing $L=\hat T+(\hat V-E)$ and expanding the squared residual norm:
\begin{equation}
\|L\Psi\|_w^2
=
\|\hat T\Psi\|_w^2 + \|(\hat V-E)\Psi\|_w^2
+2\,\Re\langle \hat T\Psi,\;(\hat V-E)\Psi\rangle_w,
\label{eq:coulomb:expand}
\end{equation}
we see that near coalescence both $\hat T\Psi$ and $\hat V_{\mathrm C}\Psi$ are large.
The exact solution relies on cancellation between these contributions
(cf.\ \eqref{eq:coulomb:EL_sing_2d_anti} and \eqref{eq:coulomb:cusp_2d_par}), which means
\eqref{eq:coulomb:expand} demands high \emph{relative} accuracy in the cusp region:
if either term is off by a small fraction, the net residual can be large in absolute magnitude.
This is a classic recipe for ill-conditioning: the loss landscape contains directions
where small parameter changes cause huge residual changes in a small region.

\paragraph{Operator-level picture (\(L^\ast L\)).}
From the same viewpoint as Appendix~\ref{app:laplacian}, strong residual training
induces matrices/operators tied to $L^\ast L$ \cite{De_Ryck_2024}.
For $L=\hat T+(\hat V-E)$, one has (formally)
\begin{equation}
L^\ast L
=
\hat T^2 + (\hat V-E)^2 + \hat T(\hat V-E) + (\hat V-E)\hat T.
\label{eq:coulomb:LstarL}
\end{equation}
Coulomb enters here in two destabilizing ways:
\begin{enumerate}
\item the positive term $(\hat V-E)^2$ involves a $1/r^2$-type singular weight;
\item the mixed terms contain derivatives of $V$ through product rules, e.g.
$\hat T(V\Psi)$ generates $\nabla V\cdot\nabla\Psi$ and $\Delta V$ contributions,
which are even more singular/distributional at coalescence.
\end{enumerate}
This formal expansion explains why the Coulomb singularity can severely inflate
the effective spectral spread of the training operator, aligning with the
operator-conditioning viewpoint \cite{De_Ryck_2024}.

\section{Barron/spectral perspective: the Coulomb singularity}
\label{app:coulomb:barron}

Coulomb forces the \emph{true} 2D solution to have cusp-like, non-smooth short-range
structure, so approximating the \emph{strong residual} requires representing
high-frequency content that is not well captured by globally smooth/low-complexity
function classes.

In spectral/Barron-style approximation theories~\cite{Barron1993-UniversalApprox}, approximation quality is governed
by spectral decay and norms that weight the Fourier tail \cite{De_Ryck_2024}.
A cusp (or a linear node times a cusp in $J$) implies slower spectral decay than a globally
smooth function.
Strong residual training then \emph{further} amplifies the high-frequency tail
because it applies $\Delta$ (and effectively $L^\ast L$), which multiplies Fourier
modes by powers of $|\mathbf k|$ (Appendix~\ref{app:laplacian}).
Hence, even if $\Psi$ can be approximated reasonably in function value, the strong
residual can remain large unless the cusp structure is represented accurately.
This is exactly the type of regularity mismatch that motivates weak/variational
formulations in numerical analysis of PINN-type methods, since weak forms reduce
the derivative order demanded of the approximation \cite{De_Ryck_2024}.

Practically, QMC wavefunctions therefore incorporate explicit short-range
correlation structure (e.g.\ cusp-enforcing Jastrow factors), precisely to remove
these singular residual pathologies and reduce the variance of local quantities
\cite{BeccaSorella2017-QMCBook,Siljamaki2003-QD-Cusps}.

\section{Implications for Stochastic Reconfiguration}
\label{app:coulomb:sr}

SR performs a natural-gradient step using a (quantum) Fisher/metric matrix
\cite{BeccaSorella2017-QMCBook,Amari1998-NaturalGradient,ParkKastoryano2020-GeometryNQS}.
In one common formulation,
\begin{equation}
S_{ij}=\mathrm{Cov}_{p_\theta}[O_i,O_j],
\qquad
g_i=2\,\mathrm{Cov}_{p_\theta}[E_L,O_i],
\qquad
S\,\delta\theta=-g,
\label{eq:coulomb:sr}
\end{equation}
with $O_i=\partial_{\theta_i}\log\Psi_\theta$.

Coulomb-induced cusp mismatch produces $E_L$ spikes of the form
\eqref{eq:coulomb:EL_mismatch_2d_anti}--\eqref{eq:coulomb:EL_mismatch_2d_par},
which has two consequences:
\begin{enumerate}
\item \textbf{Heavy-tailed Monte Carlo forces.}
The estimator of $g_i=\mathrm{Cov}[E_L,O_i]$ becomes noisy when $E_L$ has rare but
extreme spikes (short-range events). This increases variance and can destabilize
the linear solve unless one uses damping/regularisation, as standard in SR practice
\cite{BeccaSorella2017-QMCBook}.
\item \textbf{Coupling to curvature sensitivity.}
Even if SR corrects anisotropy in parameter space (natural gradient)
\cite{Amari1998-NaturalGradient,ParkKastoryano2020-GeometryNQS}, it does not remove
the underlying operator-driven stiffness: the spikes originate from imperfect
kinetic--potential cancellation enforced by Coulomb cusp physics.
\end{enumerate}

\paragraph{Summary.}
Laplacian dominance creates stiffness through $k^4$-type amplification of curvature
errors; Coulomb creates stiffness by enforcing cusp constraints (in 2D: $1$ vs $1/3$
depending on spin channel) and cancellation against the Laplacian.
Both mechanisms worsen conditioning of strong residual objectives and increase
statistical difficulty of SR through heavy-tailed local-energy fluctuations
\cite{De_Ryck_2024,BeccaSorella2017-QMCBook,Siljamaki2003-QD-Cusps}.

% =========================================================
% APPENDIX C: The Collocation--Backflow Catch-22
% =========================================================
\chapter{The collocation--backflow catch-22}
\label{app:catch22}

This appendix formalizes a structural difficulty encountered when training
backflow transformations under the \emph{strong-form} collocation objective---the
one that minimises the variance of the pointwise local energy $E_L$---and presents
extensive numerical evidence for the $N{=}6$, $\omega{=}0.5$ quantum dot.
The core observation is a \emph{catch-22}: backflow's learning signal is concentrated near
the nodal surface, but the local energy diverges there for any imperfect wavefunction, and
every sampling strategy we tried that includes those points produced gradients whose
variance we could not control. The result, across every intervention below, is an accuracy
floor at roughly $0.3\%$---set by the Jastrow alone---that is absent under VMC\,+\,SR
training of the same ansatz family. We call the floor irreducible \emph{within this
pipeline}: it is an empirical bound established by exhaustion over the interventions we
tested, not a proof that no strong-form scheme can do better.

The scope of this result needs two qualifications. First, it is a property of the
\emph{strong form} specifically: the weak-form (score-function / REINFORCE) collocation objective used
for the results in the main text never differentiates through the Laplacian and does
not inherit this pathology, which is why it trains the same ansatz to far below the
$0.3\%$ floor. The value here is therefore a diagnosis of one objective's limits, not
a measure of what Markov-chain-free training can achieve. Second, this analysis is what motivated the shift to the weak form, and is reported as a negative result and not as a headline number.

\section{Setup and notation}
\label{app:catch22:setup}

We consider a Slater--Jastrow--backflow ans\"atz of the form
\begin{equation}
\Psi_\theta(\mathbf R)
=\det\tilde S_\uparrow(\mathbf R;\theta_{\mathrm{BF}})
 \;\det\tilde S_\downarrow(\mathbf R;\theta_{\mathrm{BF}})
 \;\exp\bigl(J(\mathbf R;\theta_J)\bigr),
\label{eq:c22:ansatz}
\end{equation}
where $\tilde S_\sigma$ is the Slater matrix evaluated at backflow-shifted
coordinates $\tilde{\mathbf r}_i = \mathbf r_i + \Delta\mathbf
x_i(\mathbf R;\theta_{\mathrm{BF}})$, $J$ is a neural Jastrow correlator, and $\theta=(\theta_J,\theta_{\mathrm{BF}})$ are trainable
parameters.
A CTNN message-passing network produces the backflow shifts $\Delta\mathbf x_i$.

We train by minimising the \emph{screened collocation} objective
\begin{equation}
\mathcal L(\theta)
=\operatorname{Var}_{\mathcal S}\!\bigl[E_L(\mathbf R;\theta)\bigr],
\qquad
E_L=\frac{\hat H\Psi_\theta}{\Psi_\theta},
\label{eq:c22:loss}
\end{equation}
where the collocation set $\mathcal S$ is obtained by
importance-sampling from a Gaussian proposal~$q$
and retaining the top-$K$ points ranked by $|\Psi_\theta|^2/q$
(``screened collocation'').
This top-$K$ selection naturally suppresses near-node configurations
where $|\Psi|^2\approx 0$.


\section{The divergence of $E_L$ near nodes}
\label{app:catch22:divergence}

The local energy decomposes as
\begin{equation}
E_L
= -\frac12\frac{\nabla^2\Psi}{\Psi}+V
= -\frac12\Bigl[
  \nabla^2 J + |\nabla J|^2
  + 2\,\nabla J\cdot\frac{\nabla D}{D}
  + \frac{\nabla^2 D}{D}
\Bigr]+V,
\label{eq:c22:EL_decomp}
\end{equation}
where $D:=\det\tilde S_\uparrow\det\tilde S_\downarrow$ and
the last term $\nabla^2 D/D$ contains second derivatives of the
backflow shifts $\nabla^2\Delta\mathbf x_i$ amplified by the
inverse determinant.

\paragraph{For the exact wavefunction, $E_L$ is constant.}
If $\Psi$ is an eigenstate, $E_L=E$ everywhere. The cancellation is a property of the
\emph{full} local energy $\hat H\Psi/\Psi$ with $\Psi=D\,e^{J}$, in which the Jastrow cross
term $2\nabla J\cdot\nabla D/D$ of Eq.~\eqref{eq:c22:EL_decomp} participates alongside
$\nabla^2D/D$; we isolate the determinant terms below because they are the ones the backflow
parameters move, not because the others are absent.
For an exact eigenstate, the potentially singular contributions in the full
expression $\nabla^2(De^J)/(De^J)$ cancel at a simple node. This is not a
property of $\nabla^2D/D$ in isolation: for a generic smooth determinant with
$D\sim d$ near the nodal surface $\mathcal N$, $\nabla^2D$ may remain
$\mathcal O(1)$. The exact eigenstate satisfies
\begin{equation}
\frac{\nabla^2(De^J)}{De^J}\;\text{finite on }\mathcal N,
\label{eq:c22:exact_cancel}
\end{equation}
through cancellation among the determinant, Jastrow-cross, and remaining
Jastrow terms in Eq.~\eqref{eq:c22:EL_decomp}.

\paragraph{For an approximate wavefunction, the cancellation generically fails.}
During training, the backflow parameters $\theta_{\mathrm{BF}}$ are
imperfect: the \emph{learned} nodal surface $\mathcal N_\theta$ (where
$D(\cdot;\theta)=0$) does not coincide with the \emph{true} nodal surface.
Let $d_\theta$ denote the signed distance from the learned node and $d^*$
that from the true node.
Near a point where $d_\theta\approx 0$ but $d^*\neq 0$, the determinant
$D\sim d_\theta\to 0$ while the kinetic contributions from the Jastrow and
potential terms remain finite:
\begin{equation}
E_L(\mathbf R)
\supset
\frac{\nabla^2 D(\mathbf R;\theta)}{D(\mathbf R;\theta)}
\sim \frac{c}{d_\theta}\to\pm\infty
\qquad\text{as }d_\theta\to 0,\ d^*\neq 0,
\label{eq:c22:EL_div}
\end{equation}
for some configuration-dependent constant $c\neq 0$.
Conversely, near a point on the true node ($d^*=0$) but off the learned node
($d_\theta\neq 0$), $\Psi\sim d^*\to 0$ and the wavefunction assigns
vanishing weight.
The pathology \eqref{eq:c22:EL_div} is therefore not an artefact of the
nodal structure itself---it is a property of \emph{mismatched} nodes,
and is worst precisely when backflow still has something to learn.


\section{The geometric mismatch argument}
\label{app:catch22:mismatch}

We now locate the divergence. It is \emph{not} in the variance of the local energy, which
the calculation below shows to be finite for a simple node; it is one derivative later, in
the parameter gradient.

Consider the contribution to $\operatorname{Var}[E_L]$ from a
neighborhood $U_\varepsilon$ of the learned nodal surface:
\begin{equation}
\operatorname{Var}[E_L]
\ge
\int_{U_\varepsilon}
|\Psi(\mathbf R;\theta)|^2\,
\bigl(E_L(\mathbf R;\theta)-\bar E\bigr)^2\,
d\mathbf R.
\label{eq:c22:var_lb}
\end{equation}
When the node is ``simple'' ($D$ vanishes linearly), we have
\begin{align}
|\Psi|^2 &\sim d_\theta^2, &
E_L^2 &\sim d_\theta^{-2\alpha},
\label{eq:c22:scaling}
\end{align}
where $\alpha$ characterises the severity of the mismatch.
For the generic case where the numerator
$\nabla^2 D$ does \emph{not} vanish on $\mathcal N_\theta$
(because $\mathcal N_\theta\neq\mathcal N^*$),
one has $\alpha=1$ so that $E_L^2\sim d_\theta^{-2}$.

In $D$-dimensional configuration space, the volume element in a
tubular neighborhood of a codimension-1 surface scales as
$dV\sim d_\theta^0\,d\sigma\,dd_\theta$.
Thus the integrand in \eqref{eq:c22:var_lb} behaves as
\begin{equation}
|\Psi|^2\cdot E_L^2\cdot dV
\sim
d_\theta^2\cdot d_\theta^{-2}\cdot dd_\theta
= dd_\theta,
\label{eq:c22:integrand}
\end{equation}
which integrates to $\mathcal{O}(\varepsilon)$ and is therefore \emph{finite}: a simple node does not, by itself, make $\mathrm{Var}[E_L]$ diverge under the $|\Psi|^2$ measure. The instability enters one derivative later, in the \emph{parameter gradient}.
But the key observation is that this is the \emph{best case}.
In practice, the mismatch between $\mathcal N_\theta$ and $\mathcal N^*$
is not uniform: there exist regions where the surfaces cross at angles,
creating pockets where $|d^*|/|d_\theta|$ is large
(the true wavefunction is non-negligible while the learned determinant
is near-zero), and there the $E_L^2$ spike is not compensated by $|\Psi|^2$
suppression.

More concretely, the gradient
$\partial\mathcal L/\partial\theta_{\mathrm{BF}}$ involves terms of the form
\begin{equation}
\frac{\partial E_L}{\partial\theta_{\mathrm{BF}}}
\;\supset\;
\frac{1}{D}\frac{\partial(\nabla^2 D)}{\partial\theta_{\mathrm{BF}}}
\;-\;
\frac{\nabla^2 D}{D^2}
\frac{\partial D}{\partial\theta_{\mathrm{BF}}},
\label{eq:c22:grad_EL}
\end{equation}
both of which diverge as $D\to 0$.
Even in the marginal case where $\mathrm{Var}[E_L]$ is finite,
the \emph{gradient variance} (which involves these higher-order poles)
generically diverges, making gradient-based optimisation impossible at
points near the learned node.


\section{The catch-22}
\label{app:catch22:statement}

The preceding analysis reveals a three-sided dilemma:

\begin{enumerate}
\item \textbf{Backflow learns at the node.}
  The purpose of backflow is to reshape the nodal surface of the Slater
  determinant.
  The gradient signal $\partial E_L/\partial\theta_{\mathrm{BF}}$
  is largest where the backflow shifts actually change which side of the
  node a configuration point lies on, i.e., where $D\approx 0$.

\item \textbf{$E_L$ diverges at the node for imperfect $\theta$.}
  As shown in \S\ref{app:catch22:divergence}--\ref{app:catch22:mismatch},
  the local energy and its parameter gradient contain poles of the form
  $D^{-1}$ and $D^{-2}$ whenever the learned node does not coincide with
  the exact node. The resulting gradient variance is unbounded.

\item \textbf{Good collocation sampling must avoid the node.}
  Under screened collocation, the top-$K$ selection ranks points by
  $|\Psi|^2/q\propto D^2 e^{2J}/q$.
  Since $D^2\to 0$ at the node, near-node points are systematically
  excluded from the training batch. This is \emph{protective}---it
  prevents gradient explosion---but it also removes the points
  where backflow's learning signal lives.
\end{enumerate}

\noindent
A ``catch-22'' emerges:
\begin{quote}
\emph{Including near-node points gives divergent, noisy gradients and
backflow shrinks defensively or diverges.
Excluding them gives stable training but removes backflow's learning signal.
In either regime, collocation cannot effectively train backflow beyond an
accuracy floor set by the correlator quality alone.}
\end{quote}


\section{Why VMC\,+\,SR escapes}
\label{app:catch22:vmc}

Under VMC, the sampling distribution \emph{is} $|\Psi_\theta|^2$, which
assigns vanishing measure to the nodal surface.
Near-node configurations are strongly \emph{suppressed}---for a simple node the probability of lying within a distance $\epsilon$ scales as $\epsilon^3$---rather than strictly absent, which is enough to keep the $D^{-1}$ singularity out of the gradient estimator in practice.
However, VMC \emph{still optimises the node position} through a
subtle mechanism:

\paragraph{Indirect nodal information through bulk correlations.}
The variational energy
$\langle E_L\rangle_{|\Psi|^2}$ is sensitive to nodal position because
shifting the node changes the wavefunction \emph{everywhere}
(through normalisation and through the functional form of $D$).
The SR force vector
$g_i=\mathrm{Cov}_{|\Psi|^2}[E_L,O_i]$
(where $O_i=\partial_{\theta_i}\log\Psi$)
receives contributions from the \emph{bulk}---regions where both
$|\Psi|^2$ and $E_L$ are well-behaved---that are
\emph{correlated with} nodal position but do not require evaluating $E_L$
at the node.

\paragraph{SR amplifies the signal.}
The Fisher metric $S_{ij}=\mathrm{Cov}[O_i,O_j]$ captures how sensitive
$\log\Psi$ is to each parameter direction.
By solving $(S+\varepsilon I)\,\delta\theta=-g$, SR
amplifies exactly those subtle bulk signals that encode nodal information,
compensating for the fact that the node itself has zero weight under
$|\Psi|^2$.

\paragraph{Core asymmetry.}
The asymmetry needs stating carefully, because a power count that is too quick proves too
much. Under $|\Psi|^2$ the SR force $g_i=\mathrm{Cov}[E_L,O_i]$ has a finite \emph{mean}:
both $E_L$ and $O_i=\partial_{\theta_i}\log\Psi$ carry a $d_\theta^{-1}$ pole near a
mismatched node, their product goes as $d_\theta^{-2}$, and the measure supplies
$|\Psi|^2\sim d_\theta^{2}$, so the integral converges. Its \emph{variance} does not, by the
same counting---so VMC is not immune to near-node pathology either, and in practice this is
what the damping and trust region of \S\ref{sec:sr-stage} are for.

The real distinction is one derivative higher. VMC and SR never differentiate $E_L$: the
local energy enters only as a multiplicand in a covariance, so the $D^{-2}$ pole of
Eq.~\eqref{eq:c22:grad_EL}---which comes from $\partial_\theta$ acting on $\nabla^2D/D$---never
appears. Strong-form collocation differentiates $E_L$ by construction, and inherits it.
That is why the same node is survivable in one scheme and not in the other, and it is the
argument \S\ref{sec:colloc-why-reinforce} makes for the weak form as well.
None of the sampling, loss, or preconditioning strategies we tested
(\S\ref{app:catch22:experiments}) overcame this asymmetry within the strong-form
collocation framework; we present it as a structural obstruction of the tested pipeline
rather than a theorem that no collocation method can ever avoid it.


\section{Experimental evidence}
\label{app:catch22:experiments}

We present a systematic study for $N{=}6$ electrons in a 2D harmonic trap
with $\omega{=}0.5$ (DMC reference: $E_{\mathrm{DMC}}=11.78484(6)$
\cite{Pederiva2000-QD-DMC}).
The ans\"atz uses a neural Jastrow correlator (${\sim}15{,}000$ parameters)
and a CTNN backflow network (${\sim}8{,}600$ parameters) with
$\mathrm{bf\_scale}=0.7$.
All energies are evaluated via MCMC with 15\,000 samples after
burn-in.

\subsection{Baseline: correlator alone vs.\ correlator\,+\,backflow}
\label{app:catch22:baseline}

\begin{table}[htbp]
\centering
\caption[Baseline strong-form collocation results]{Baseline collocation results for $N{=}6$, $\omega{=}0.5$.
``Correlator only'' trains without backflow; ``joint'' trains both
networks from scratch for 300 epochs with $\mathrm{bf\_scale}=0.7$.
The backflow provides only a marginal improvement of ${\sim}0.1$ percentage
points.}
\label{tab:c22:baseline}
\begin{tabular}{l c c c}
\toprule
Configuration & $E$ (Ha) & $\sigma_E$ & err (\%) \\
\midrule
Correlator only (300 ep) & 11.834 & 0.003 & 0.42 \\
Correlator\,+\,BF, joint (300 ep) & 11.823 & 0.003 & 0.33 \\
Correlator\,+\,BF, separate opt & 11.822 & 0.003 & 0.32 \\
\midrule
DMC reference & 11.78484 & --- & 0.00 \\
\bottomrule
\end{tabular}
\end{table}

The backflow reduces the error from $0.42\%$ to $0.32$--$0.33\%$, far short
of the ${<}0.01\%$ reached by VMC\,+\,SR on the same ansatz \emph{family}
(Table~\ref{tab:energies}). The two are not matched in size: the networks here total
${\sim}24$k parameters against ${\sim}236$k for the production model, so the comparison
bounds what this objective achieves at this capacity, not isolating the objective alone. The floor is nonetheless a property of the objective, because it does not move under
any of the interventions below, including ones that add capacity.

\subsection{Sampling strategies}
\label{app:catch22:sampling}

We tested multiple strategies to bring near-node information into the
collocation batch, hoping to give the backflow a stronger learning signal.

\begin{table}[htbp]
\centering
\caption[Effect of sampling strategy on collocation accuracy]{Effect of sampling strategy on collocation accuracy.
All experiments use 50-epoch fine-tuning from the best
collocation checkpoint ($0.35\%$), except ``hard-smooth (scratch)''
which trains 300 epochs from random initialisation.}
\label{tab:c22:sampling}
\begin{tabular}{l l c c}
\toprule
Method & Description & err (\%) & Outcome \\
\midrule
Top-$K$ baseline & standard screened collocation & 0.35 & stable \\
Hard mining + Huber & top-$K$ of residuals & 0.41 & neutral \\
Hard mining + smoothness & + $\|\nabla^2\Delta x\|^2$ penalty & 0.38 & slight help \\
SIR (single proposal) & $|\Psi|^2$-reweighted resampling & 0.47 & worse \\
SIR (mixture) & Gaussian mixture proposal & 0.44 & worse \\
Gumbel top-$K$~\cite{JangGuPoole2017-GumbelSoftmax} & stochastic top-$K$ selection & 5.00 & catastrophic \\
\bottomrule
\end{tabular}
\end{table}

Every strategy that includes more near-node points
(SIR, Gumbel) produced \emph{worse} results---consistent with the
catch-22: the additional points carry divergent $E_L$ values that
corrupt the gradient.
Gumbel top-$K$ (which replaces deterministic selection with stochastic
Gumbel-softmax) was catastrophic, with early stopping triggered by
rapid energy divergence.

\subsection{Loss modifications}
\label{app:catch22:loss}

\begin{table}[htbp]
\centering
\caption[Effect of loss-function modifications]{Effect of loss function modifications.
The hybrid loss adds a mean-energy (variational) term;
the det filter discards points with small $|\det\tilde S|$.}
\label{tab:c22:loss}
\begin{tabular}{l l c c}
\toprule
Method & Description & err (\%) & Outcome \\
\midrule
Variance loss (baseline) & Huber on $E_L-\bar E$ & 0.33 & best collocation \\
Hybrid anneal & $\beta\cdot\text{Var}+(1{-}\beta)\cdot\text{Mean}$,
  $\beta\!:\!0.7{\to}0.15$ & 2.59 & catastrophic \\
Gumbel + hybrid & stochastic top-$K$ + hybrid & 5.00 & catastrophic \\
Hybrid + det filter & $\beta{=}0.2$ fixed + $|\det|$ filter & 24.75 & broken \\
\bottomrule
\end{tabular}
\end{table}

The mean-energy gradient $\nabla_\theta\langle E_L\rangle_{\mathcal S}$
under biased collocation sampling does \emph{not} equal the VMC gradient---the
sampling distribution is not $|\Psi|^2$---so the variational principle
does not hold, and the backflow drifts to large,
unphysical shifts ($|\Delta\mathbf x|\approx 0.83$, versus ${\sim}0.23$
under variance-only training).
The variance loss's implicit constraint on backflow magnitude is
\emph{protective}, not pathological.

\subsection{Gradient regularisation}
\label{app:catch22:regularisation}

We attacked the divergence sources in $E_L$ directly:

\begin{table}[htbp]
\centering
\caption[Regularisation of the $E_L$ singularity]{Regularisation of the $E_L$ singularity (50-epoch
fine-tune from $0.35\%$ checkpoint). Each modification smooths or clips
one of the divergent contributions to $E_L$.}
\label{tab:c22:reg}
\begin{tabular}{l l c c}
\toprule
Method & Modification & err (\%) & Outcome \\
\midrule
Soft Coulomb & $1/r\to 1/\sqrt{r^2+\varepsilon^2}$,
  $\varepsilon{=}0.2$ & 1.41 & much worse \\
Regularised det & $\log|D|\to\operatorname{logaddexp}(\log|D|,\,
  \log\delta)$ & 0.49 & worse \\
$T$ clipping & $\max(|T|,T_{\max})$, $T_{\max}{=}30$ & 0.43 & neutral \\
\bottomrule
\end{tabular}
\end{table}

Soft Coulomb trains on wrong physics ($V$ is no longer the true potential).
Regularising the determinant distorts the $\log\Psi$ landscape.
Kinetic energy clipping is neutral because the clipped points are exactly
those that carried the (noisy) backflow signal.
None of the regularisations improve upon the baseline.

\subsection{Gradient preconditioning}
\label{app:catch22:precondition}

\begin{table}[htbp]
\centering
\caption[Gradient preconditioning experiments]{Gradient preconditioning experiments (300-epoch scratch training
unless noted). K-FAC~\cite{MartensGrosse2015-KFAC} uses block-diagonal Kronecker-factored
Fisher on the backflow network only.}
\label{tab:c22:kfac}
\begin{tabular}{l c c}
\toprule
Method & err (\%) & Outcome \\
\midrule
Vanilla Adam & 0.33 & baseline \\
K-FAC on CTNN & 0.34 & no improvement \\
K-FAC + smoothness & 0.51 & worse \\
Variance-minimising pretrain & 0.54 & worse \\
Adam fine-tune (30 ep, lr$=$1e-4) & 0.34 & no improvement \\
Adam burst (15 ep, lr$=$3e-4) & 0.56 & worse \\
\bottomrule
\end{tabular}
\end{table}

K-FAC provides a block-diagonal approximation to the Fisher matrix
and should, in principle, help with the ill-conditioned parameter-space
geometry.
In practice, it does not: the conditioning problem is
\emph{operator-level} (Appendices~\ref{app:laplacian}--\ref{app:coulomb}),
not merely parameter-space anisotropy, and K-FAC cannot address the
fundamental divergence in $E_L$.

\subsection{Summary of all interventions}
\label{app:catch22:summary}

\begin{table}[htbp]
\centering
\caption[All strong-form collocation interventions]{Complete summary of all collocation-based backflow experiments
for $N{=}6$, $\omega{=}0.5$.
No intervention breaks the ${\sim}0.32\%$ accuracy floor.
The final row is the same ansatz \emph{family} under VMC\,+\,SR at production width
(Table~\ref{tab:energies}), an order of magnitude larger in parameter count; it is included
to show the size of the gap, not as a matched comparison.}
\label{tab:c22:summary}
\begin{tabular}{l c l}
\toprule
Experiment & err (\%) & Category \\
\midrule
Correlator\,+\,BF, sep opt (best) & 0.32 & baseline \\
Correlator\,+\,BF, joint 300 ep & 0.33 & baseline \\
K-FAC on CTNN & 0.34 & preconditioner \\
Adam fine-tune 30 ep & 0.34 & optimiser \\
Top-$K$ baseline (retrained) & 0.35 & sampling \\
Hard mining + smoothness & 0.38 & sampling \\
Hard mining + Huber & 0.41 & sampling \\
Correlator only & 0.42 & no backflow \\
$T$ clipping & 0.43 & regularisation \\
SIR (mixture) & 0.44 & sampling \\
SIR (single) & 0.47 & sampling \\
Regularised det & 0.49 & regularisation \\
K-FAC + smoothness & 0.51 & preconditioner \\
Var-min pretrain & 0.54 & pretrain \\
Adam burst lr$=$3e-4 & 0.56 & optimiser \\
Soft Coulomb & 1.41 & regularisation \\
Hybrid loss (annealed) & 2.59 & loss \\
Gumbel top-$K$ + hybrid & 5.00 & sampling + loss \\
Hybrid + det filter & 24.75 & loss + filter \\
\midrule
\textbf{Same ansatz, VMC\,+\,SR} & $<$\textbf{0.01} &
  \textbf{variational} \\
\bottomrule
\end{tabular}
\end{table}

The pattern is striking:
\begin{itemize}
\item Every intervention that \emph{includes} more near-node points
  (SIR, Gumbel) \emph{increases} the error---backflow shrinks
  defensively under gradient noise or diverges entirely.
\item Every intervention that \emph{changes} the loss to tolerate
  near-node contributions (hybrid mean-energy, det filter) produces
  catastrophic results because the collocation sampling is not $|\Psi|^2$
  and the variational principle does not apply.
\item The smoothness penalty $\lambda\|\nabla^2\Delta\mathbf x\|^2$
  helps modestly (from $0.59\%$ to $0.38\%$ without it vs.\ with it)
  by damping the $\nabla^2 D/D$ poles, but cannot eliminate the
  fundamental divergence.
\item Ironically, the only method that trains backflow effectively is
  MCMC-based VMC\,+\,SR---which sidesteps the entire collocation
  framework.
\end{itemize}


\section{Connection to Appendices~\ref{app:laplacian} and \ref{app:coulomb}}
\label{app:catch22:connection}

The catch-22 is a specific manifestation of the general conditioning
difficulties analysed in Appendices~\ref{app:laplacian}--\ref{app:coulomb}:

\paragraph{Laplacian stiffness (Appendix~\ref{app:laplacian}).}
The $k^4$-type amplification of curvature errors
(\S\ref{app:laplacian:bilap}) is compounded by backflow, because
$\nabla^2(\det\tilde S)$ introduces second derivatives of
$\Delta\mathbf x$ that are not present in a fixed-node ansatz.
These additional high-frequency modes inflate the spectral spread
of the training operator $L^\ast L$.

\paragraph{Coulomb singularity (Appendix~\ref{app:coulomb}).}
The $1/r$ Coulomb potential (\S\ref{app:coulomb:setup}) adds a
singular multiplication operator to $\hat H$.
Near both the coalescence point ($r_{ij}\to 0$) and the node
($D\to 0$), $E_L$ contains products of these two singularities,
and the cusp cancellation (\S\ref{app:coulomb:cusp}) must be
maintained simultaneously with the nodal cancellation
\eqref{eq:c22:exact_cancel}.
In the strong residual squared $\|LE_L\|^2$, these compound
singularities create an extremely stiff landscape.

\paragraph{The backflow-specific pathology.}
What distinguishes the backflow catch-22 from the general conditioning
problem is the \emph{self-referential} nature of the singularity:
the backflow parameters $\theta_{\mathrm{BF}}$ determine \emph{where}
the node is, so updating $\theta_{\mathrm{BF}}$ moves the singularity
in configuration space.
Under the collocation loss, the gradient pushes $\theta_{\mathrm{BF}}$
in a direction determined by $E_L$ values at the current collocation
points, but the resulting node shift invalidates those very points.
VMC+SR avoids this because it samples from the \emph{current}
$|\Psi_\theta|^2$ at every step, so the sampling distribution
co-adapts with the node position.


% =========================================================
% APPENDIX D: Additional collocation results
% =========================================================
\chapter{Additional collocation results}
\label{app:postcatch22}

This appendix records the collocation results that Chapter~\ref{ch:kernel} refers to but does
not tabulate: the $N{=}20$ energies, and the variants that were tried and failed. Both bear on
the same question---what limits the Markov-chain-free route---from the side of what does not
work.

\section{What was varied}
\label{app:postcatch22:programme}

Four axes were explored. The \emph{sampling}: Gaussian and Gaussian-mixture proposals scaled
by $1/\sqrt\omega$, oversample-then-resample with ESS diagnostics, ESS-adaptive oversampling,
weight tempering and log-weight clipping. The \emph{objective}: finite-difference collocation
residuals, a hybrid weak form, and pure REINFORCE, which became the default, with the
Laplacian kept in the reward path but never differentiated. The \emph{preconditioner}: Adam,
diagonal-Fisher natural gradient, and full SR in Woodbury and CG modes with damping anneals
and trust-region caps. And the \emph{training style}: single sweeps, wave-based cascades with
checkpoint inheritance, and asynchronous campaigns relaunching from the best available
checkpoint.

Two results held across the search. Optimisation quality dominated architectural novelty:
moderately sized backflow$+$Jastrow models with stable training beat more ambitious variants
with weaker optimisation geometry. And at low $\omega$, regularising the resampling
distribution (ESS control, top-mass reduction) helped where hard cusp gating did not---a
controlled $2{\times}2{\times}3$ matrix over gating, resampling regularisation and three
seeds showed no measurable gain from gating and a consistent gain from regularisation.

Warm-starting along the confinement cascade
$\omega:1.0\!\to\!0.5\!\to\!0.1\!\to\!0.01\!\to\!0.001$ was necessary throughout; cold starts
at low $\omega$ did not converge. Representative cascade results are $N{=}6$ at
$\omega{=}1.0$ ($E=20.16167\pm0.00181$, $+0.012\%$), $\omega{=}0.5$
($E=11.78466\pm0.00199$, $-0.002\%$) and $\omega{=}0.1$ ($E=3.56525\pm0.00046$, $+0.32\%$),
and $N{=}12$ at $\omega{=}1.0$ ($E=65.70699\pm0.00595$, $+0.010\%$). For $N{=}2$ the route
covers the whole confinement range: $+0.012\%$, $+0.023\%$, $-0.004\%$, $+0.002\%$ and
$+0.154\%$ at $\omega=1,0.5,0.1,0.01,0.001$, zero-variance extrapolated.

\section{The $N=20$ collocation energies}
\label{app:postcatch22:frontier}

At $N{=}20$ the backflow's edge-feature tensor, which scales as
$\mathcal{O}(N^2 h_{\mathrm{bf}})$, forces the hidden dimension down to
$h_{\mathrm{bf}}{=}48$--$64$ to fit in memory. The best backflow run at $\omega{=}1.0$ then
reached only ${\sim}{+}18\%$ error, while a Jastrow-only ansatz---freed of that tensor, and
so able to afford a larger collocation budget and higher oversampling---reached $+1.32\%$
($E=157.934\pm0.018$). The more expressive architecture is beaten by the simpler one, which
is consistent with the memory-forced capacity cut degrading gradient quality more than its
extra flexibility repays.

Polishing the Jastrow checkpoints at reduced learning rate improved the $\omega{=}1.0$ error
from ${\sim}{+}2.6\%$ to $+1.32\%$ over ${\sim}400$ epochs, and the $\omega{=}0.5$ error from
${\sim}{+}7.0\%$ to $+2.74\%$ ($E=96.447\pm0.015$); the best $\omega{=}0.1$ result was
$+5.53\%$ ($E=31.637\pm0.008$). The gap to the sub-percent accuracy reached at $N\in\{6,12\}$
tracks the capacity cut and the gradient noise it brings. A more memory-frugal backflow is
the direct next intervention; these experiments do not isolate the collocation objective as
the source of the remaining gap.

\section{Negative results worth recording}
\label{app:postcatch22:summary}

Scalar final energy is not sufficient for method development on this route: ESS, top-mass
concentration, rollback counts and the probe-to-final gap are what distinguish a run that
converged from one that stopped moving. Several intuitively attractive interventions
delivered nothing in the tested range---hard cusp gating most notably---and recording that
was necessary to keep the search from circling. Stricter SR trust constraints, heavier
damping, ESS-adaptive sampling and rollback controls improved operational stability without
closing the accuracy gap at the hardest targets. What eventually closed it at low $\omega$
and $N{=}6$ was not a stabilisation measure but a change of proposal
(\S\ref{sec:kernel-q3-frontier}).


\chapter{The learned representation of one trained solution}
\label{app:representation}

This appendix records what the internal representation of \emph{one} high-quality trained
solution looks like. It is descriptive, it is measured on a single training run, and---as
\S\ref{sec:synthesis} sets out---most of the quantities in it are component-level and
therefore gauge-like: the correlator and the backflow are partial substitutes, so how the
labour is divided between them records the route the training took as much as the state it
reached. Nothing in the architectural argument of Chapter~\ref{ch:kernel} rests on anything
here. It is reported because the regularities are striking and because they are ones the
network was free \emph{not} to adopt, not because they establish a property of the
architecture.

Three findings survive that qualification well enough to be worth stating. Across
$N\in\{2,6,12,20\}$ and $\omega\in\{1,0.5,0.1,10^{-2},10^{-3}\}$ the correlator feature
space stays on one to three effective dimensions and never expands, even where the
real-space density develops elaborate shell structure; what changes across the crossover is
which branch the leading axis is built from, not how many axes there are. The displacement
field is the opposite kind of object, occupying nearly all the spatial directions available
to it. And the coupling between the two changes character near $\omega\approx0.1$, the same
place the energy partition crosses over (\S\ref{sec:wigner-energy}).

Every diagnostic below is evaluated on fresh samples drawn from $|\Psi_\theta|^2$, with the
estimators and stability protocol of \S\ref{sec:repr-methods}.

%% ──────────────────────────────────────────────────────────────────────
\section{Correlator geometry: a low-rank representation that reorganises}
\label{sec:repr-geom}

We characterise the correlator feature space $Z = \overline\phi\|\overline\psi\|\mathbf{g}$
by its entropy-based effective rank $r_{\mathrm{eff}}(Z)$, the correlation
$|\rho(\mathrm{head},\mathrm{PC1})|$ between the head output and its first principal component,
and the power of PC1 decomposed across architectural branches.
Table~\ref{tab:correlator-geom} reports these quantities across the grid.

\begin{table}[t]
\centering
\caption[Correlator feature geometry]{Correlator geometry (CTNN).
$r_{\mathrm{eff}}$: entropy-based effective rank of $Z$.
$|\rho|$: absolute correlation between head output and PC1 reconstruction.
$\phi/\psi/\mathrm{g}$: fraction of PC1 variance from the orbital, pair, and auxiliary branches.
The $\psi$-dominated entries ($\psi > 0.5$) are shown in bold.}
\label{tab:correlator-geom}
\small
\setlength{\tabcolsep}{4pt}
\begin{tabular}{rr cccccc}
\toprule
$N$ & $\omega$ & $r_{\mathrm{eff}}$ & $|\rho|$ & $\phi$ & $\psi$ & $\mathrm{g}$ \\
\midrule
 2 & 0.001 & 1.00 & 1.00 & 0.00 & \textbf{1.00} & 0.00 \\
 2 & 0.01  & 1.04 & 0.66 & 0.00 & \textbf{1.00} & 0.00 \\
 2 & 0.1   & 1.31 & 0.99 & 0.00 & \textbf{0.53} & 0.46 \\
 2 & 0.5   & 1.25 & 1.00 & 0.00 & \textbf{0.58} & 0.42 \\
 2 & 1.0   & 1.53 & 0.98 & 0.00 & 0.35 & \textbf{0.65} \\
\midrule
 6 & 0.001 & 1.52 & 0.97 & 0.01 & 0.12 & \textbf{0.87} \\
 6 & 0.01  & 1.38 & 0.99 & 0.05 & 0.44 & \textbf{0.50} \\
 6 & 0.1   & 2.73 & 0.94 & \textbf{0.66} & 0.21 & 0.14 \\
 6 & 0.5   & 2.69 & 0.98 & 0.15 & \textbf{0.63} & 0.22 \\
 6 & 1.0   & 2.69 & 0.99 & 0.03 & \textbf{0.62} & 0.35 \\
\midrule
12 & 0.001 & 1.66 & 0.96 & \textbf{0.54} & 0.13 & 0.33 \\
12 & 0.01  & 1.54 & 0.97 & 0.43 & 0.09 & \textbf{0.48} \\
12 & 0.1   & 3.08 & 0.99 & 0.12 & \textbf{0.63} & 0.26 \\
12 & 0.5   & 2.32 & 0.99 & 0.12 & \textbf{0.69} & 0.19 \\
12 & 1.0   & 1.76 & 0.99 & 0.08 & \textbf{0.73} & 0.19 \\
\midrule
20 & 0.001 & 2.59 & 0.99 & 0.18 & \textbf{0.74} & 0.08 \\
20 & 0.01  & 1.09 & 0.79 & 0.04 & \textbf{0.94} & 0.01 \\
20 & 0.1   & 2.20 & 0.97 & 0.36 & 0.44 & 0.20 \\
20 & 0.5   & 2.30 & 1.00 & 0.17 & \textbf{0.70} & 0.13 \\
20 & 1.0   & 2.32 & 0.99 & 0.18 & \textbf{0.71} & 0.11 \\
\bottomrule
\end{tabular}
\end{table}

\paragraph{The manifold is low-dimensional everywhere.}
The effective rank stays below $3.1$ across the entire grid,
and the head output is almost perfectly reconstructed from PC1 alone in all but
one analysed cell ($|\rho|\geq0.94$ except $N{=}20$, $\omega{=}0.01$, where
$|\rho|=0.79$).
This holds even at strong correlation where the real-space density develops pronounced shell
structure---the correlator does not respond by expanding its latent dimensionality.
The modest peak at $N{=}12$, $\omega{=}0.1$ ($r_{\mathrm{eff}} = 3.08$) coincides with the
intermediate-correlation regime where neither the shell picture nor the crystalline picture provides
a clean single-mode description.

\paragraph{A $\phi\leftrightarrow\psi$ transition in the leading axis.}
What varies across regimes is \emph{which branch dominates PC1}.
At moderate-to-weak correlation ($\omega \geq 0.1$, $N\geq 6$),
PC1 is predominantly $\psi$-loaded: the pair branch carries 60--73\% of the variance,
meaning the correlator's primary degree of freedom tracks pair correlations.
At strong correlation the loading shifts toward $\phi$ (orbital) or $\mathbf{g}$
(auxiliary) branches for $N{=}6$ and $N{=}12$. The $N{=}20$ run does not follow
the same branch-switching pattern: its pair branch remains dominant at both weakest
confinements. Given that these are single-run, component-level diagnostics, we do
not interpret this discrepancy as a size-dependent physical transition.
The transition is sharpest at $N{=}12$, where $\psi$ goes from 0.09 at $\omega{=}0.01$
to 0.63 at $\omega{=}0.1$.

These loadings track the physics.
When interactions are weak relative to the trap ($\omega \gtrsim 0.1$),
the non-interacting shell structure provides a good zeroth-order description,
and the correlator's primary role is refining pair correlations on top of that.
When interactions dominate ($\omega \lesssim 0.01$),
the leading correlator coordinate shifts away from pair features toward the
single-particle and auxiliary branches. The network reorganises the features through
which it represents the state, and that reorganisation---not any growth in
dimension---is how it adapts.
Nothing in the architecture imposes this behaviour: the $N{=}6$ and $N{=}12$
regime changes appear as a rotation of a low-dimensional axis across branches rather
than as growth in latent dimensionality.

%% ──────────────────────────────────────────────────────────────────────
\section{Backflow: a high-rank field with a low-rank footprint}
\label{sec:repr-backflow}

Table~\ref{tab:backflow-geom} characterises the message-passing displacement field.
We report the normalised displacement magnitude $\|\Delta\mathbf{x}\|/\ell$,
the effective rank of $\Delta\mathbf{x}$ (viewed as a $B\times 2N$ matrix of flattened displacements),
and the CKA similarity between the correlator features with and without backflow.

\begin{table}[t]
\centering
\caption[Backflow displacement geometry]{CTNN backflow geometry.
$\|\Delta\mathbf{x}\|/\ell$: mean displacement magnitude normalised by harmonic length.
Rank: effective rank of the displacement covariance.
$2N{-}2$: maximum possible rank after centre-of-mass projection.
CKA: linear centred kernel alignment between correlator features with and without backflow
(1.0 = identical feature manifolds).}
\label{tab:backflow-geom}
\small
\setlength{\tabcolsep}{5pt}
\begin{tabular}{rr cccc}
\toprule
$N$ & $\omega$ & $\|\Delta\mathbf{x}\|/\ell$ & Rank & $2N{-}2$ & CKA \\
\midrule
 2 & 0.001 & 0.003  & 1.9  &  2 & 1.000 \\
 2 & 1.0   & 0.029  & 2.0  &  2 & 1.000 \\
\midrule
 6 & 0.001 & 0.017  & 10.0 & 10 & 0.999 \\
 6 & 0.01  & 0.035  & 9.9  & 10 & 0.998 \\
 6 & 0.1   & 0.096  & 10.0 & 10 & 0.995 \\
 6 & 0.5   & 0.093  & 10.0 & 10 & 0.991 \\
 6 & 1.0   & 0.089  & 10.0 & 10 & 0.979 \\
\midrule
12 & 0.001 & 0.017  & 21.9 & 22 & 1.000 \\
12 & 0.01  & 0.030  & 21.9 & 22 & 0.998 \\
12 & 0.1   & 0.096  & 21.9 & 22 & 0.993 \\
12 & 0.5   & 0.096  & 21.9 & 22 & 0.993 \\
12 & 1.0   & 0.091  & 21.9 & 22 & 0.995 \\
\midrule
20 & 0.001 & 0.023  & 37.7 & 38 & 0.999 \\
20 & 0.01  & 0.046  & 37.4 & 38 & 1.000 \\
20 & 0.1   & 0.075  & 37.8 & 38 & 0.997 \\
20 & 0.5   & 0.099  & 37.6 & 38 & 0.990 \\
20 & 1.0   & 0.099  & 37.6 & 38 & 0.991 \\
\bottomrule
\end{tabular}
\end{table}

\paragraph{Displacement magnitude scales with $\omega$ but saturates near $\|\Delta\mathbf{x}\|/\ell \sim 0.1$.}
At strong correlation ($\omega \leq 0.01$), displacements are small relative to the harmonic length,
$\|\Delta\mathbf{x}\|/\ell \sim 0.02$--$0.05$.
As the trap tightens and the system enters the moderate-to-weak correlation regime,
the magnitude grows and saturates at $\sim$\,8--10\% of $\ell$ for $\omega \gtrsim 0.1$.
The message-passing backflow never produces large displacements---its $\tanh$ output bound and copresheaf message-passing structure
act as implicit regularisers, in contrast to the conventional backflow, which can exhibit
displacements comparable to the system scale in poorly converged regimes
(e.g.\ $\|\Delta\mathbf{x}\|/\ell \sim 0.95$ for $N{=}6$, $\omega{=}0.001$ with the conventional network).

\paragraph{Full-rank displacements, near-unity CKA.}
The effective rank of the displacement field is consistently $\approx 2N - 2$,
saturating the available spatial degrees of freedom modulo centre-of-mass conservation
(confirmed by COM shifts of order $10^{-14}$--$10^{-17}$).
The backflow transformation is not a low-rank collective mode
such as a breathing or rigid rotation; the displacement field spans essentially all COM-free spatial directions, so each
particle receives a correction
conditioned on the full configuration.
The field stays near that ceiling across the whole grid, including the deepest Wigner
points, where one might have expected the correction to become collective. Whether that
is a property of \emph{this} architecture or of trained backflows generally is not
something the present measurement can decide; Chapter~\ref{ch:kernel} settles it by
holding everything else fixed (\S\ref{sec:kernel-q1b}).

Yet the CKA between correlator features with and without backflow exceeds 0.99 in nearly all cases
(and exceeds 0.979 everywhere).
The correlator manifold is almost unperturbed by the displacement field.
The correlator's feature geometry is therefore almost unaffected by the presence of the
displacement field, which is a statement about that geometry and not about the subspace the
backflow occupies. Whether the two components are genuinely complementary, rather than
merely uncorrelated on this measure, is tested by ablation in \S\ref{sec:kernel-q1b}.

%% ──────────────────────────────────────────────────────────────────────
\section{Correlator--backflow coupling: from independence to cooperation}
\label{sec:repr-coupling}

The previous subsection established that backflow leaves the correlator manifold intact.
But are the correlator and the backflow \emph{functionally} coupled---does the displacement
depend on what the correlator computes?
We measure this through three diagnostics:
(i) the mean cosine similarity $\langle\cos(\nabla_{\mathbf{x}} f,\,\Delta\mathbf{x})\rangle$
between the correlator's gradient in configuration space and the displacement field
(alignment implies the backflow pushes along directions the correlator is sensitive to);
(ii) the fraction of configurations where $\Delta\mathbf{x}$ is aligned with $\nabla f$;
and (iii) $|\delta f|/\sigma_f$, where $\delta f = f(\tilde R+\Delta\mathbf x)-f(\tilde R)$
is the change in the correlator output when it is evaluated on the backflow-shifted
coordinates, in units of the correlator's own standard deviation $\sigma_f$. Two
cautions attach to this last quantity. It is a \emph{counterfactual sensitivity}: in the
production ansatz the correlator sees the unshifted $\tilde R$ (only the determinant sees
$\tilde R+\Delta\mathbf x$; \S\ref{subsec:backflow-method}), so $\delta f$ measures how
much the correlator \emph{would} respond to the displacement, not an actual term in the
wavefunction. And it is a typical relative change, not a variance decomposition:
$|\delta f|/\sigma_f=0.7$ means the counterfactual change is $0.7$ standard deviations,
not that $70\%$ of the variance is attributable to backflow.
Table~\ref{tab:pinn-ctnn-coupling} reports these alongside the fraction of particles pushed
toward their nearest neighbour.

\begin{table}[t]
\centering
\caption[Correlator--backflow coupling diagnostics]{Correlator--backflow coupling diagnostics (message-passing ansatz).
$\langle\cos\rangle$: mean cosine between $\nabla f$ and $\Delta\mathbf{x}$.
Aligned: fraction of configs with $\nabla f \cdot \Delta\mathbf{x} > 0$.
$|\delta f|/\sigma_f$: backflow-induced change in the correlator output, relative to its standard deviation.
$\rightarrow$NN: fraction of particles displaced toward their nearest neighbour ($0.5 =$ isotropic).
Only $N\geq 6$ shown; $N{=}2$ displacements are near-zero for most $\omega$.}
\label{tab:pinn-ctnn-coupling}
\small
\setlength{\tabcolsep}{4.5pt}
\begin{tabular}{rr cccc}
\toprule
$N$ & $\omega$ & $\langle\cos\rangle$ & Aligned & $|\delta f|/\sigma_f$ & $\rightarrow$NN \\
\midrule
 6 & 0.001 &  0.14 & 0.63 & 0.03 & 0.43 \\
 6 & 0.01  &  0.18 & 0.77 & 0.05 & 0.28 \\
 6 & 0.1   &  0.49 & 0.99 & 0.25 & 0.16 \\
 6 & 0.5   &  0.62 & 1.00 & 0.33 & 0.19 \\
 6 & 1.0   &  0.77 & 1.00 & 0.41 & 0.18 \\
\midrule
12 & 0.001 & $-$0.48 & 0.01 & 0.03 & 0.61 \\
12 & 0.01  &  0.22 & 0.82 & 0.05 & 0.29 \\
12 & 0.1   &  0.56 & 1.00 & 0.39 & 0.12 \\
12 & 0.5   &  0.65 & 1.00 & 0.49 & 0.20 \\
12 & 1.0   &  0.75 & 1.00 & 0.55 & 0.21 \\
\midrule
20 & 0.001 &  0.32 & 0.97 & 0.06 & 0.44 \\
20 & 0.01  &  0.22 & 0.89 & 0.07 & 0.45 \\
20 & 0.1   &  0.39 & 1.00 & 0.25 & 0.17 \\
20 & 0.5   &  0.65 & 1.00 & 0.58 & 0.22 \\
20 & 1.0   &  0.78 & 1.00 & 0.69 & 0.25 \\
\bottomrule
\end{tabular}
\end{table}

\paragraph{A transition at $\omega\approx 0.1$.}
The data reveal a sharp regime change.
At $\omega \leq 0.01$ the coupling is weak and inconsistent in sign: $\langle\cos\rangle$
runs from $-0.48$ to $0.32$ across the six cells, the alignment fraction is erratic, and
backflow barely affects the correlator output ($|\delta f|/\sigma_f < 0.07$ everywhere). The
$N{=}12$, $\omega{=}0.001$ cell is strongly \emph{anti}-aligned
($\langle\cos\rangle=-0.48$, alignment fraction $0.01$), which we record without an
account of it: at that confinement the displacement field is small enough
(Table~\ref{tab:backflow-geom}) that its direction relative to the correlator gradient
carries little energy either way, but we cannot say why the sign is systematic rather than
scattered.
In this regime the backflow operates in a subspace largely orthogonal to the correlator's gradient---it
modifies the wavefunction in directions that $W_\theta$ is insensitive to,
consistent with its role being to adjust the Slater determinant's orbital positions rather than
the correlator.

At $\omega \geq 0.1$, the picture changes qualitatively.
The gradient cosine jumps to 0.5--0.8, the alignment fraction reaches 1.0 for $N\geq 6$,
and the counterfactual output change reaches $|\delta f|/\sigma_f = 0.25$--$0.69$ (i.e.\ up to $\sim\!0.7$ standard deviations of the correlator).
Here the backflow acts as a cooperative refinement: it pushes configurations along
directions the correlator is maximally sensitive to, systematically correcting the correlator's
residual errors.
The transition is sharpest between $\omega=0.01$ and $\omega=0.1$,
coinciding with the crossover in $\Gamma = V_{\mathrm{int}}/T$ from $\Gamma > 10$
(interaction-dominated) to $\Gamma \sim 3$--$6$ (moderate correlation).

\paragraph{Directional structure: correlation-hole deepening.}
The nearest-neighbour direction fraction provides a complementary signature.
At $\omega \geq 0.1$ for $N\geq 6$, only 12--25\% of particles are displaced toward their
nearest neighbour.
The backflow preferentially pushes particles \emph{apart}, actively deepening the correlation hole.
This directional bias is strongest at $N{=}12$, $\omega{=}0.1$ ($\rightarrow$NN\,$=0.12$),
exactly where the correlator effective rank peaks and the $\phi\leftrightarrow\psi$ transition occurs.
At strong correlation ($\omega\leq 0.01$, large $N$), the fraction rises toward 0.5
(isotropic), because the large equilibrium inter-particle separations leave little room for
further correlation-hole deepening.

%% ──────────────────────────────────────────────────────────────────────

\section{Limits of component-level interpretation}
\label{sec:repr-synthesis}

Taken together these measurements describe two components of very different character. The
correlator occupies one to three directions, whose leading axis tracks the dominant global
coordinate of the cloud and changes identity rather than growing in number as the physics
changes. The displacement field is high-rank and configuration-dependent, and its presence
leaves the correlator's feature geometry almost unchanged (CKA $>0.98$). The coupling
between them is weak and inconsistent in sign at strong correlation, and strong and
systematically aligned at moderate coupling, where the displacements follow the
correlator's gradient; the change is sharp and coincides with the crossover in $\Gamma$.

These measurements describe one trained solution and do not by themselves establish which
properties are architecture-specific. A different training route could divide the same
labour differently and arrive at the same energy, which for two partial substitutes is what
one should expect. Deciding the question requires holding everything else fixed and varying
only the architecture.


\chapter{Structural diagnostics in full}
\label{app:structural-tables}

Chapter~\ref{ch:results} reports the two radial quantities the argument turns
on---$r_{\rm mode}$ and the width-to-scale ratio $\gamma$. The complete diagnostics are
collected here: the mean radius, width, FWHM and decile quantiles for $N{=}2$, and the
ring-angle spread and pair Lindemann ratio for $N{=}6$. Estimators and stability
protocol are defined in \S\ref{sec:analysis-struct}. All lengths are in Bohr; angles in
radians.

\section{Fitted energy and radius exponents}
\label{app:energy-exponents}

Section~\ref{sec:wigner-energy} states two conclusions from the energy partition---that
interactions come to dominate, and that the trap--interaction balance approaches its
classical form---and defers the fitted exponents that support them to here.

Log--log fits $E(\omega)\propto\omega^{\alpha}$ over the grid
$\omega\in\{10^{-3},10^{-2},10^{-1},0.5,1\}$ give a kinetic exponent
$\alpha_T\simeq1.00$ at every $N$, an interaction exponent
$\alpha_{V_{\rm int}}\simeq0.66$, and trap exponents
$\alpha_{V_{\rm trap}}\simeq0.81,\,0.73,\,0.72$ for $N=2,6,12$. The kinetic energy therefore
scales essentially linearly with $\omega$ while both potential terms grow sublinearly with
nearly $N$-independent exponents; the ratio $\Gamma\propto\omega^{\alpha_{V_{\rm
int}}-\alpha_T}\approx\omega^{-0.34}$ predicts a factor $10^{3\times0.34}\approx10.5$ across
three decades, against the measured $9.7$, $10.6$ and $11.1$.

From $V_{\rm trap}=\tfrac12\omega^2\sum_i\langle r_i^2\rangle$ the rms radius is
$r_{\rm rms}=\sqrt{2V_{\rm trap}/(N\omega^2)}$, and fitting
$r_{\rm rms}\propto\omega^{\alpha_r}$ gives $\alpha_r\simeq-0.60$ to $-0.64$, equivalently
$\langle r^2\rangle\propto\omega^{-\beta}$ with $\beta\simeq1.19$--$1.28$. The cloud
therefore expands faster than the non-interacting $\omega^{-1}$ law for $\langle r^2\rangle$.

These are averaged quantities and should not be confused with the shell-resolved radius of
\S\ref{app:structural-tables:radius}, which is the one compared against the classical Wigner
prefactor; the two differ because $r_{\rm rms}$ and $r_{\rm mode}$ mix the central electron
of a $(1,n)$ topology with the ring, and the shell-resolved radius does not.

\begin{table}[H]
\centering
\caption{Radial diagnostics for $N{=}2$ across $\omega$ (Bohr units).}
\label{tab:two_e_summary}
\begin{tabular}{lcccccccc}
\toprule
$\omega$
& $r_{\rm mode}$
& $\langle r\rangle$
& $\sigma_r$
& FWHM
& $q_{10}$
& $q_{50}$
& $q_{90}$
& $\gamma$ \\
\midrule
1.00  & 1.435  & 1.632  & 0.705  & 1.743  & 0.735  & 1.581  & 2.581  & 0.491 \\
0.50  & 2.303  & 2.475  & 1.007  & 2.466  & 1.209  & 2.409  & 3.829  & 0.437 \\
0.10  & 3.601  & 3.830  & 1.728  & 4.449  & 1.640  & 3.714  & 6.127  & 0.480 \\
0.01  & 14.191 & 14.593 & 5.615  & 13.602 & 7.390  & 14.322 & 22.038 & 0.396 \\
0.001 & 64.292 & 63.953 & 18.479 & 43.977 & 40.392 & 63.814 & 87.714 & 0.287 \\
\bottomrule
\end{tabular}
\end{table}

\begin{table}[H]
\centering
\caption{Radial and angular diagnostics for $N{=}6$ across $\omega$
(Bohr; angles in radians).}
\label{tab:N6}
\begin{tabular}{c|ccccc}
\toprule
$\omega$ & $r_{\rm mode}$ & $\sigma_r$ & $\gamma$
& $\sigma(\text{ring angles})$ & $\gamma_{r_{ij}}$ \\
\midrule
1.00  & 2.0907   & 0.9249  & 0.4424 & 0.7046 & 0.3538 \\
0.50  & 3.0702   & 1.3648  & 0.4445 & 0.6943 & 0.3397 \\
0.10  & 8.1686   & 3.5025  & 0.4288 & 0.6842 & 0.3149 \\
0.01  & 33.1538  & 13.9094 & 0.4195 & 0.6497 & 0.2889 \\
0.001 & 113.5542 & 47.8076 & 0.4210 & 0.6368 & 0.2664 \\
\bottomrule
\end{tabular}
\end{table}

\section{Shell-resolved structure across the whole grid}
\label{app:structural-tables:shells}

Chapter~\ref{ch:results} quotes shell occupancies, radii, bond order and Lindemann
numbers for individual $(N,\omega)$ cells in running prose. Table~\ref{tab:shell-summary}
collects them for the whole grid from a single analysis pass, so that the trends can be read off instead of assembled from the text. Shell radii are given in trap units
$\tilde r=\sqrt{\omega}\,r$; multiply by $\ell=\omega^{-1/2}$ for Bohr.

\paragraph{Reading $|\Phi_n|$ against its null.}
One column of that table needs a baseline to be interpretable, and its absence would
make the bond-order numbers read far more favourably than they should. For $n$ particles
placed at \emph{independent uniform} angles on a ring, $\Phi_n=n^{-1}\sum_k e^{in\phi_k}$
is a two-dimensional random walk of $n$ unit steps, so
\begin{equation}
  \big\langle|\Phi_n|\big\rangle_{\rm random}=\frac{\sqrt\pi}{2\sqrt n}\approx\frac{0.886}{\sqrt n},
  \qquad
  \big\langle|\Phi_n|^2\big\rangle_{\rm random}=\frac1n .
  \label{eq:phi-null}
\end{equation}
A finite $|\Phi_n|$ is therefore \emph{not} by itself evidence of angular order: a
five-particle ring with no order at all returns $|\Phi_5|\approx0.40$. The ratio
$|\Phi_n|/\langle|\Phi_n|\rangle_{\rm random}$ is the quantity that carries information, and
Table~\ref{tab:shell-summary} reports it.

Read that way the result is sharper than the raw values suggest, and it needs stating with
the right threshold. From $\omega=1$ down to $\omega=0.1$ the ratio is $1.00$--$1.02$ at
every particle number and on every ring. At $\omega=0.01$ it lifts slightly, to $1.10$ at
$N{=}6$, while $N{=}12$ and $N{=}20$ remain at $1.03$ and $1.01$. Only at $\omega=10^{-3}$
does it become large: $1.65$ on the $N{=}6$ ring, $1.19$ and $1.56$ on the two rings of the
$N{=}12$ $(3,9)$ molecule, and $1.19$ and $1.47$ on the outer two rings of the $N{=}20$
$(1,7,12)$ molecule.

Whether the small excesses are \emph{resolved} is a separate question from whether they are
\emph{large}, and with these sample sizes the two come apart. Two caveats attach to the null
before it is used quantitatively. The Rayleigh form is the large-$n$ limit of the resultant
of $n$ unit steps; for the small rings here ($n=3$ to $19$) it is an approximation, accurate
to well within the effect at $10^{-3}$ but of the same order as the few-percent excesses at
$\omega\ge0.1$. And the frames come from a Markov chain, so a naive standard error over $K$
raw frames understates the uncertainty, whereas the block bootstrap of
\S\ref{subsec:repr-uncertainty} is used for the reported errors elsewhere. Taking those at
face value, the excess at $N{=}6$, $\omega=0.1$ is of order tens of block standard errors and
still only $2\%$.

The conclusion that survives both caveats is the one that matters: statistical detectability
is not what separates the regimes here, magnitude is. On that criterion the ordering is flat
through $\omega=0.1$, weakly present at $\omega=0.01$ and only at $N{=}6$, and unambiguous at
$10^{-3}$. A cleaner version of this comparison would generate the null numerically at each
actual occupancy $n$, through the same block-bootstrap pipeline, instead of relying on the asymptotic form; that is a small addition we did not make.

The Lindemann numbers of the same table fall monotonically throughout, so radial stiffening
precedes $n$-fold angular ordering rather than accompanying it---the separation between
radial and angular melting anticipated in \S\ref{sec:wigner-markers-theory}. Note that the
values here are for the \emph{modal} topology over all frames; the sector-conditioned values
quoted in \S\ref{sec:wigner-angular} are computed on the frames assigned to a given sector
and run slightly higher.

\begin{table}[H]
\centering
\caption[Shell-resolved diagnostics across the whole grid]{Shell-resolved diagnostics at the modal topology, all frames, for every
$(N,\omega)$. Occupancies are $(n_{\rm in},\dots,n_{\rm out})$; $\tilde r_s$ are mean shell
radii in trap units; $w_s=\sigma(r_s)/\langle r_s\rangle$ are per-shell radial relative
widths; $|\Phi_{n_s}|$
is the rotation-registered bond order on shell $s$; and the final column is $|\Phi_{n_s}|$
divided by its random-angle expectation, Eq.~\eqref{eq:phi-null}. A value of $1.00$ in that column means the ring is
equal to the Rayleigh null approximation within quoted precision; see the text for
why resolvability and magnitude come apart at these sample sizes. Single-particle inner shells are omitted from the
angular columns, where $|\Phi_1|=1$ trivially.}
\label{tab:shell-summary}
\small
\setlength{\tabcolsep}{5pt}
\begin{tabular}{ccl cc cc}
\toprule
$N$ & $\omega$ & occ. & $\tilde r_s$ & $w_s$ & $|\Phi_{n_s}|$ & ratio to null \\
\midrule
 6 & $1.0$   & $(1,5)$     & $0.70,\,1.66$ & $0.43,\,0.32$ & $0.397$ & $1.00$ \\
 6 & $0.5$   & $(1,5)$     & $0.76,\,1.78$ & $0.42,\,0.31$ & $0.397$ & $1.00$ \\
 6 & $0.1$   & $(1,5)$     & $0.93,\,2.15$ & $0.41,\,0.28$ & $0.405$ & $1.02$ \\
 6 & $0.01$  & $(1,5)$     & $1.33,\,2.96$ & $0.42,\,0.22$ & $0.434$ & $1.10$ \\
 6 & $0.001$ & $(1,5)$     & $1.12,\,4.29$ & $0.63,\,0.15$ & $\mathbf{0.654}$ & $\mathbf{1.65}$ \\
\midrule
12 & $1.0$   & $(1,11)$    & $0.55,\,2.00$ & $0.47,\,0.34$ & $0.268$ & $1.00$ \\
12 & $0.5$   & $(1,11)$    & $0.60,\,2.18$ & $0.47,\,0.33$ & $0.268$ & $1.00$ \\
12 & $0.1$   & $(1,11)$    & $0.74,\,2.68$ & $0.45,\,0.31$ & $0.270$ & $1.01$ \\
12 & $0.01$  & $(1,11)$    & $1.04,\,3.77$ & $0.41,\,0.28$ & $0.275$ & $1.03$ \\
12 & $0.001$ & $(3,9)$     & $2.52,\,6.04$ & $0.61,\,0.24$ & $\mathbf{0.611,\,0.462}$ & $\mathbf{1.19,\,1.56}$ \\
\midrule
20 & $1.0$   & $(1,19)$    & $0.55,\,2.32$ & $0.44,\,0.36$ & $0.204$ & $1.00$ \\
20 & $0.5$   & $(1,19)$    & $0.60,\,2.52$ & $0.44,\,0.35$ & $0.203$ & $1.00$ \\
20 & $0.1$   & $(1,19)$    & $0.75,\,3.14$ & $0.42,\,0.34$ & $0.204$ & $1.00$ \\
20 & $0.01$  & $(1,19)$    & $1.05,\,4.47$ & $0.49,\,0.32$ & $0.205$ & $1.01$ \\
20 & $0.001$ & $(1,7,12)$  & $1.27,\,4.19,\,7.83$ & $0.47,\,0.24,\,0.11$ & $\mathbf{0.398,\,0.375}$ & $\mathbf{1.19,\,1.47}$ \\
\bottomrule
\end{tabular}
\end{table}

\section{Shell-model fit to the pair distribution}
\label{app:structural-tables:mixfit}

The reconstruction quoted in \S\ref{sec:wigner-angular} tests whether the measured pair
distribution is explained by the detected shell geometry. Two versions of that test exist
and they are not equally informative.

The weaker version decomposes the pairs of two-shell frames into inner--inner,
inner--outer and outer--outer histograms and recombines them with the combinatorial weights
of the ideal topology. Because the sector histograms and the total are accumulated from the
same frames, this recombination is close to exact by construction---cosine similarities
exceed $0.99999$ at every $(N,\omega)$---and it establishes only that the sector
bookkeeping is self-consistent.

The stronger version is a \emph{generative} one and is the test reported here. Shell radii
and widths are fitted on one half of the frames; each shell is then modelled as a regular
$n_s$-gon carrying a per-frame global rotation and a fitted angular jitter; and the
resulting pair-distance density is compared with the empirical density on the \emph{held-out}
half. Against it we run an angle-randomised null with the same fitted radial distributions
but uniform angles, which isolates what polygonal order contributes over and above shell
radii alone. Table~\ref{tab:shell-model-fit} reports both.

\begin{figure}[htbp]
  \centering
  \includegraphics[width=0.8\textwidth]{mixfit_N12_w0.001_nin3.pdf}
  \caption[Shell-resolved reconstruction of $g(r)$, $N{=}12$]{Shell-resolved reconstruction of $g(r)$ at $\omega=10^{-3}$ for
  $N{=}12$.  The global pair distribution is reproduced by a convex mixture
  of inner--inner (II), inner--outer (IO) and outer--outer (OO) histograms
  with weights set by the observed $(n_{\rm in},n_{\rm out})$ statistics.}
  \label{fig:N12_w0001_decomposition}
\end{figure}

The polygon model reproduces the held-out pair distribution to cosine similarity
$0.987$--$0.999$ across the whole grid. The null reaches $0.936$--$0.996$, and the gap
between them, $\Delta\cos$, grows monotonically as the trap weakens at every particle
number---from $0.012$ to $0.058$ at $N{=}6$, from $0.006$ to $0.026$ at $N{=}12$, and from
$0.003$ to $0.010$ at $N{=}20$. Shell radii alone account for most of the pair structure at
strong confinement; angular order accounts for a growing remainder as the Wigner molecule
forms. This is the same transition the bond-order ratios of
Table~\ref{tab:shell-summary} record, measured on a different quantity and on held-out data.

\begin{table}[H]
\centering
\caption[Shell-model fit to the pair-distance distribution]{Shell-model fit to the pair-distance distribution. The empirical
$p_{\rm test}(r_{ij})$ on held-out frames is compared with (i) a \emph{polygon} model, in
which each detected shell is a regular $n_s$-gon with a per-frame global rotation and fitted
angular jitter, and (ii) a \emph{uniform-angle} null with the same fitted radial
distributions and no polygonal order. $\Delta\cos=\cos_{\rm poly}-\cos_{\rm uni}$ is what
angular order buys; $L^2_{\rm poly}$ is the polygon model's residual.}
\label{tab:shell-model-fit}
\small
\setlength{\tabcolsep}{6pt}
\begin{tabular}{ccl cccc}
\toprule
$N$ & $\omega$ & occ. & $\cos_{\rm poly}$ & $\cos_{\rm uni}$ & $\Delta\cos$ & $L^2_{\rm poly}$ \\
\midrule
 6 & $1.0$   & $(1,5)$    & $0.9988$ & $0.9866$ & $0.0122$ & $4.69\times10^{-3}$ \\
 6 & $0.5$   & $(1,5)$    & $0.9986$ & $0.9836$ & $0.0150$ & $5.09\times10^{-3}$ \\
 6 & $0.1$   & $(1,5)$    & $0.9974$ & $0.9741$ & $0.0234$ & $6.83\times10^{-3}$ \\
 6 & $0.01$  & $(1,5)$    & $0.9922$ & $0.9536$ & $0.0386$ & $1.14\times10^{-2}$ \\
 6 & $0.001$ & $(1,5)$    & $0.9942$ & $0.9358$ & $\mathbf{0.0585}$ & $1.04\times10^{-2}$ \\
\midrule
12 & $1.0$   & $(1,11)$   & $0.9993$ & $0.9937$ & $0.0056$ & $3.36\times10^{-3}$ \\
12 & $0.5$   & $(1,11)$   & $0.9990$ & $0.9927$ & $0.0063$ & $4.01\times10^{-3}$ \\
12 & $0.1$   & $(1,11)$   & $0.9980$ & $0.9896$ & $0.0085$ & $5.43\times10^{-3}$ \\
12 & $0.01$  & $(1,11)$   & $0.9929$ & $0.9805$ & $0.0124$ & $9.92\times10^{-3}$ \\
12 & $0.001$ & $(3,9)$    & $0.9868$ & $0.9604$ & $\mathbf{0.0264}$ & $1.34\times10^{-2}$ \\
\midrule
20 & $1.0$   & $(1,19)$   & $0.9992$ & $0.9964$ & $0.0028$ & $3.43\times10^{-3}$ \\
20 & $0.5$   & $(1,19)$   & $0.9990$ & $0.9958$ & $0.0032$ & $3.78\times10^{-3}$ \\
20 & $0.1$   & $(1,19)$   & $0.9982$ & $0.9942$ & $0.0040$ & $4.95\times10^{-3}$ \\
20 & $0.01$  & $(1,19)$   & $0.9947$ & $0.9900$ & $0.0047$ & $8.21\times10^{-3}$ \\
20 & $0.001$ & $(1,7,12)$ & $0.9915$ & $0.9820$ & $\mathbf{0.0095}$ & $1.02\times10^{-2}$ \\
\bottomrule
\end{tabular}
\end{table}

\section{The measured shell radius against the classical prediction}
\label{app:structural-tables:radius}

The classical minimum-energy configuration of six charges in a parabolic trap places five
of them on a ring of radius $1.334\,\omega^{-2/3}$
\cite{schweigert1994spectralpropertieschargedparticles,Kong_2002}. Nothing in the ansatz,
the training objective or the sampler knows that number, so the fitted outer-shell radius of
Table~\ref{tab:shell-summary} is an external check on the whole structural picture---and a
sharper one than the exponent fits of \S\ref{sec:wigner-radial}, because it tests a
prefactor as well as a power.

\begin{table}[H]
\centering
\caption[Measured shell radius against the classical prediction]{Fitted outer-shell radius of the $N{=}6$ $(1,5)$ molecule against the classical
Wigner ring $r_{\rm cl}=1.334\,\omega^{-2/3}$. Radii in Bohr.}
\label{tab:classical-radius}
\begin{tabular}{ccccc}
\toprule
$\omega$ & $r_{\rm shell}$ & $r_{\rm cl}$ & ratio \\
\midrule
$1.0$   & $1.66$   & $1.33$   & $1.247$ \\
$0.5$   & $2.52$   & $2.12$   & $1.188$ \\
$0.1$   & $6.80$   & $6.19$   & $1.098$ \\
$0.01$  & $29.65$  & $28.74$  & $1.032$ \\
$0.001$ & $135.67$ & $133.40$ & $1.017$ \\
\bottomrule
\end{tabular}
\end{table}

The ratio falls monotonically toward unity, reaching $1.017$ at $\omega=10^{-3}$: the
quantum ring sits $1.7\%$ outside the classical one at the weakest trap, and the excess
grows smoothly with confinement as zero-point motion pushes the shell outward. Fitting
$r_{\rm shell}\propto\omega^{\alpha}$ gives $\alpha=-0.637$ over the full grid and $-0.650$
over the weakest three confinements, against the classical $-2/3$.

Two remarks on why this quantity, rather than the $r_{\rm mode}$ of
Table~\ref{tab:radial-compact}, is the right one to compare. First, $r_{\rm mode}$ is the
mode of the \emph{single-particle} radial density, which for a $(1,5)$ topology mixes the
central electron with the ring and is therefore not the ring radius; fitting it gives
$\alpha\approx-0.58$ at $N{=}6$ and $-0.54$ at $N{=}2$ over the full grid, and approaches
$-2/3$ only on the weak-trap half. Second, the shell-resolved radius is the quantity the
physics-informed proposal of \S\ref{sec:kernel-q3-frontier} actually needs, so the agreement
recorded here is what licenses building a sampler out of it.
